\section{Introduction}
\label{sec:introduction}

Let $\DD=\{z\in\CC:|z|<1\}$, and write $dA=dx\,dy$ for planar
Lebesgue measure. A domain is a connected open subset of $\CC$; a
univalent map is an injective holomorphic map. Brennan's conjecture asks
whether, for every conformal bijection $\varphi:W\to\DD$ from a simply
connected plane domain, the derivative satisfies
\[
 \int_W|\varphi'(z)|^s\dd A(z)<\infty
 \qquad\left(\frac43<s<4\right).
\]
There is no regularity assumption on the boundary and no boundedness
assumption on $W$. The problem measures how strongly a conformal map can
compress area near an arbitrary boundary. The conjecture dates to
Brennan's 1978 work \cite{Brennan1978}; the domain and integral-means
formulations are recorded in \cite[Chapter~3]{Bertilsson1999}.

The inverse map $f=\varphi^{-1}:\DD\to W$ puts all source domains on
the same footing. Conformal change of variables gives
\begin{equation}\label{eq:inverse-area}
 \int_W|\varphi'|^s\dd A
 =\int_{\DD}|f'|^{2-s}\dd A.
\end{equation}
Thus the conjectured interval becomes $-2<t<2/3$, with $t=2-s$.
Positive powers detect large derivatives, whereas the difficult negative
powers detect small derivatives. The inverse-square power $t=-2$ is
the limiting case relevant to the upper endpoint $s=4$.

To state its radial form, let
\[
 \SSS=\{f:\DD\to\CC\text{ univalent}:f(0)=0,\ f'(0)=1\}
\]
be the normalized schlicht class. For $f\in\SSS$, $t\in\mathbb R$,
and $0<r<1$, define
\[
 M_t[f'](r)=\frac1{2\pi}\int_{-\pi}^{\pi}
                  |f'(re^{i\theta})|^t\dd\theta,
\]
and put
\begin{align*}
 \beta_f(t)&=\limsup_{r\uparrow1}
       \frac{\log M_t[f'](r)}{\log(1/(1-r))},\\
 B_{\SSS}(t)&=\sup_{f\in\SSS}\beta_f(t).
\end{align*}
Equivalently, $\beta_f(t)$ is the infimum of the real $b$ for which
$M_t[f'](r)=O((1-r)^{-b})$ as $r\uparrow1$; the implicit constant may
depend on $f$. In particular, the supremum defining $B_{\SSS}$ is taken after forming
each map's limiting exponent. We prove the following uniform estimate
directly, and then derive the area assertion.

\begin{theorem}[Brennan's conjecture and inverse-square means]
\label{thm:main}
For every $\varepsilon>0$ there is $C_\varepsilon<\infty$ such that
\begin{equation}\label{eq:main-uniform}
 M_{-2}[f'](r)\le C_\varepsilon(1-r)^{-1-\varepsilon}
 \qquad\left(f\in\SSS,\ \frac12\le r<1\right).
\end{equation}
Moreover, $B_{\SSS}(-2)=1$.

Let $W\subset\CC$ be any simply connected domain whose boundary in
the Riemann sphere contains at least two points. Every conformal
bijection $\varphi:W\to\DD$ satisfies
\[
 \int_W|\varphi'(z)|^s\dd A(z)<\infty
 \qquad\left(\frac43<s<4\right).
\]
Equivalently, every univalent map $f:\DD\to\CC$ satisfies
\[
 \int_{\DD}|f'(w)|^t\dd A(w)<\infty
 \qquad\left(-2<t<\frac23\right).
\]
\end{theorem}

The domain correspondence, including the boundary condition, is
justified in Section~\ref{sec:integrability}. That section also derives
$B_{\SSS}(t)=|t|-1$ for every $t\le-2$; the corresponding uniform
radial bounds retain an arbitrary positive loss. The radial exponent $1$
is compatible with divergence of the area integral at $t=-2$.
The other area endpoint also fails for the same classical example.

\begin{corollary}[Sharpness of the open range]\label{cor:sharp-endpoints}
The Koebe function $f(z)=z/(1-z)^2$ belongs to $\SSS$ and satisfies
\[
 \int_{\DD}|f'|^{-2}\dd A=\infty,
 \qquad
 \int_{\DD}|f'|^{2/3}\dd A=\infty.
\]
Its inverse therefore fails area integrability at both $s=4$ and $s=4/3$.
\end{corollary}

\subsection{The problem and its development}

Brennan's 1978 work \cite{Brennan1978,Brennan1978LOMI} improved the
previously known upper range beyond $s=3$, as recorded in
Bertilsson's historical account \cite[Chapter~3, pp.~43--44]{Bertilsson1999}.
The standard formulation asks for the full interval $4/3<s<4$ on every
simply connected plane domain admitting a conformal bijection to the disk.
Neither boundedness nor boundary regularity is imposed. The classical
slit-plane obstructions at both endpoints are verified in
Section~\ref{sec:integrability}.

Subsequent progress connected the area question to finer estimates for
derivative integral means. Pommerenke developed this approach in
\cite{Pommerenke1985IntegralMeans,Pommerenke1985IntegralMeansII}.
Carleson and Makarov related negative integral means to the distribution
of large harmonic measure and proved the linear spectral tail
$B_{\SSS}(t)=|t|-1$ for all sufficiently negative $t$
\cite[\S1.4]{CarlesonMakarov1994}. The proposed threshold $t=-2$
is the one associated with Brennan's conjecture
\cite{HedenmalmShimorin2005}. The statement for sufficiently negative parameters alone
does not reach that endpoint.

Bertilsson studied coefficient estimates for negative powers of
derivatives of univalent functions \cite{Bertilsson1998}. His thesis
reports the area range $4/3<s<3.421$ \cite{Bertilsson1999} and
establishes equivalences between the spectral endpoint and
uniform integral-means formulations
\cite[Theorem~3.7, especially parts (a) and (f), pp.~44--45]{Bertilsson1999}.
Thus uniform estimates have an established place in the formulation of
Brennan's problem. In the present argument, the common constant in
\eqref{eq:main-uniform} arises directly from an operator on the entire
compact normalized class.

Hedenmalm and Shimorin combined area inequalities with weighted Bergman
space methods to improve general bounds for the integral-means spectrum.
Their published table gives $B_{\SSS}(-2)\le1.218$
\cite[Table~1, p.~385]{HedenmalmShimorin2005}.
Sola extended the Bergman-space computation
\cite{Sola2006}, reporting the improved estimate
$B_{\SSS}(-2)\le1.206$ in his thesis \cite[p.~6]{Sola2007Thesis}.
These estimates illustrate the remaining gap between general analytic
bounds and the conjectured exponent $1$.

Recent works treat the conjectured integrability range under additional
hypotheses. Aleman and Kouroupis prove it for members of continuous
semigroups of conformal disk self-maps and for Koenigs maps of nontrivial
such semigroups
\cite[Theorem~1.1 and Proposition~2.2]{AlemanKouroupis2025}.
Zheng treats simply connected polynomial basins at infinity for
$P(z)=z^m+a_0$, $m\ge3$, and for monic degree-$m$ polynomials whose
coefficients of $z,\ldots,z^{m-1}$ have modulus smaller than a
degree-dependent threshold; these basins
are considered in the sphere and contain infinity
\cite[Theorems~1.3--1.4]{Zheng2026Basins}.
Jin computes the real integral-means spectra of normalized rational
functions univalent in the disk,
with inverse-square supremum $1$ over that class
\cite[Theorem~6.20 and Corollaries~6.22--6.23]{Jin2025IntegralMeans}.
His work on Schwarzian multiplication operators recalls Shimorin's
sufficient multiplier-norm criterion for the inverse-square endpoint
of an individual normalized univalent function
\cite[Proposition~1.4(2)]{Jin2026Multipliers}.
Theorem~\ref{thm:main} concerns the full class of univalent disk maps.

\subsection{Analytic and geometric ingredients}

Normalized changes of base point, distortion estimates, and compactness
of the schlicht class are classical tools of geometric function theory
\cite[Chapters~1--2]{Duren1983}. A related probabilistic method for
exterior integral means appears in the conformal snowflakes of Beliaev
and Smirnov. Their positive integral operator gives lower bounds for an
averaged exponent defined from expected derivative integral means.
For each real $t$ and $\varepsilon>0$, they construct a random snowflake
whose averaged exponent exceeds $B(t)-\varepsilon$, where $B$ is their
universal spectrum defined using exterior Riemann maps
\cite[Lemma~3.4 and Theorem~4.1]{BeliaevSmirnov2010}.
Here the operator acts on all normalized half-plane maps, and
hypothetical excess growth produces a weighted affine law. The
operator is defined in Section~\ref{sec:preliminaries};
Section~\ref{sec:weighted-law} constructs its eigenmeasure and proves
the averaging identity.

The critical-point comparison is a planar index argument. Its
boundary calculation uses Hopf's secant proof of the tangent-turning
theorem \cite[\S2, pp.~53--55]{Hopf1935}, while the regular-target
step is an elementary plane case of Sard's theorem
\cite[Theorem~4.1]{Sard1942}.
We keep the index of a possible pole separate from those of the finite
critical points. The pairing is defined by ranks in ordered lists;
its measurability is proved through continuous local inverse branches,
and its equivariance follows from the ordering itself.

The subsequent area comparison is related in form to equivariant
mass transport, including the general framework for nonunimodular
group actions developed by Gentner and Last \cite{GentnerLast2011}.
We use a direct calculation with local changes of variables and
explicit expanding boxes. The affine covariance, exceptional-target
transfer, and boundary-cost estimates needed here are all proved
within the paper.

An independent companion applies a related affine-law and critical-point
comparison to prove a uniform inverse-first-power bound with exponent
strictly below $1/4$ for normalized univalent disk maps, disproving
Kraetzer's bounded-spectrum prediction at $p=-1$
\cite[Theorem~1.1]{OpenAIStrictInverseFirst2026}.
The proof here does not use that strict estimate; its H\"older
interpolation from \eqref{eq:main-uniform} gives exponent $1/2$ with
an arbitrary positive loss at $t=-1$.


\subsection{The obstacle and the proof strategy}

The classical distortion theorem bounds $|f'(z)|^{-1}$ by a constant
times $(1-|z|)^{-1}$ for $f\in\SSS$. Squaring this estimate gives
radial growth with exponent $2$, while the required exponent is $1$.
The missing information concerns how often very small derivatives can
occur around a circle. Our proof encodes hypothetical excess growth in
a probability law on conformal maps. That law imposes an algebraic sign
on an expected Jacobian; a geometric pairing of critical points imposes
the opposite sign.

We work on the upper half-plane $\HH=\{x+iy:y>0\}$ with the compact
class $\KK$ of univalent maps $F$ normalized by $F(i)=0$ and $F'(i)=1$.
Write $Q_F=1/F'$. Changing the base point to $z=x+iy$ gives
\[
 T_zF(w)=\frac{F(x+yw)-F(x+iy)}{yF'(x+iy)}.
\]
Each $T_zF$ again belongs to $\KK$. A positive operator averages the
two changes of base point $z=\pm1/2+i/2$, with weights $|Q_F(z)|^2$.
Its iterates sample equally spaced points above a fixed interval at
dyadic heights. If $\rho$ denotes the exponential growth rate of their
operator norms, Section~\ref{sec:preliminaries} shows that $\rho\le2$
implies the uniform disk estimate in Theorem~\ref{thm:main}.

Suppose instead that $\rho>2$, and put $\beta=\log_2\rho>1$.
Section~\ref{sec:weighted-law} constructs a Borel probability law $P$
on $\KK$ with the following property: for every $x+iy\in\HH$ and
every nonnegative Borel function $\Phi$ on $\KK$,
\[
 \mathbb E_P\bigl[|Q_F(x+iy)|^2\Phi(T_{x+iy}F)\bigr]
       =y^{-\beta}\mathbb E_P\Phi(F).
\]
We first construct a positive eigenmeasure, then average horizontal
translations and logarithmic scales. The resulting identity holds at
every affine change of base point. Differentiating the constant-test
identity determines the moments of the logarithmic derivative that
will control the Jacobian.

For $\xi\in\CC$, consider the target potential
\[
 V_\xi^F(x+iy)=\frac{y^k}{|F(x+iy)-\xi|},
 \qquad k=\frac{\beta+3}{4}>1.
\]
At the unique possible preimage of $\xi$ we give it value $+\infty$.
A critical-point count is useful only if points above a fixed positive
level cannot escape to the real axis or to infinity. Section~\ref{sec:escape}
proves precisely this compactness for almost every target and almost
every sampled map. The main boundary estimate controls the total length
of a horizontal line mapped into a compact target set.

The critical points of $V_\xi^F$ solve
\[
 G_F(z)=\xi,\qquad G_F(x+iy)=F(x+iy)-\frac{iy}{k}F'(x+iy).
\]
Section~\ref{sec:critical-pairing} shows that the sign of the normalized
real Jacobian of $G_F$ distinguishes saddles from maxima when $\xi$
is a regular value. For these targets, an index count on compact
superlevel boundaries gives at least as many maxima as saddles above
each positive level. Ordering the two lists by their values of $V_\xi^F$
therefore pairs each saddle with a distinct maximum of at least the
same value. We prove that this assignment is Borel and respects every
affine change of source coordinates. These properties are essential
when the pointwise pairing is integrated against $P$.

Section~\ref{sec:transport} performs that integration. Change of
variables through $G_F$ converts a weighted source area into the target
weight $(V_\xi^F/k)^4$. The pairing increases this weight. Affine
covariance then turns the expected source density into $dx\,dy/y$;
expanding boxes have asymptotically negligible boundary enlargement.
This proves that the expected normalized Jacobian is nonpositive,
contradicting its strictly positive value from the logarithmic-derivative
moments. A separate application of the same area formula handles the
exceptional targets at the distinguished point $i$, whose target
$G_F(i)=-i/k$ is fixed by normalization.

Finally, Section~\ref{sec:integrability} obtains the negative area
exponents $-2<t<0$ by H\"older's inequality and treats the positive exponents
with classical growth estimates and a weighted area argument. It
also proves the equivalence with the endpoint spectrum and the Koebe
sharpness assertions. Thus the analytic law, the compactness of
superlevels, and the measurable pairing each supply a distinct input
to the same weighted area contradiction.
